\begin{statement}{P3K.25}{Refuted; corrected scalar-fan computation}
\par\noindent\textbf{Audit correction, 17 November 2026.} The asserted all-shear equality is false. In the explicitly constructed primitive OS45 family of P3K.10/P3K.22, the actual scalar loop has J=1 and reference physical class 3 by K68--69. P3K.64 gives class 2 on \(T_2\), whose directions are \(e_2\),(1,2),(1,1). The direct initial/reflected C2 equality remains the narrow comparison proved in P3K.62. The outer C1 determinant has the previously omitted winding -2J. The earlier statement and argument below are preserved as the located prior claim.
\medskip

For each of the ten scalar-parameter families of P3K.22, retain its
actual period group, cap gluing, physical delta circle, and positive
scalar-parameter loop. Replace the vertical-edge cusp fan with edge
directions
\[
 e_2,\quad e_1,\quad e_1-e_2
\]
by any of the fans with edge directions
\[
 e_2,\quad e_1-j e_2,\quad e_1-(j+1)e_2,\qquad j\in\mathbb Z.
                                                        \tag{P3K.25.1}
\]
These are the actual periodic small resolutions of the same
height-one model with the original period group. Their return
indices and oriented return classes are equal to those in P3K.22:
the values are respectively \(-3,-6,-9,-14\), and their residues
are \(25,22,19,14\), in the profile groups displayed there.
The conclusion includes the OS45 construction of P3K.10.

The present proof uses the actual nonfixed-edge Cox bordism and
charged caps of P3K.17. It does not identify the repaired local
source page with the old graph trace or discard its additional
line class.
\end{statement}
\begin{proof}[Preserved earlier argument; scope corrected above]
\par\medskip\noindent\textbf{The literal square chart and the unchanged period group.}\enspace

First take the horizontal edge from \(0\) to \(e_1\). Its two
adjacent triangles have the other vertices \(e_2,e_1-e_2\).
Order the four rays, including their height-one coordinate, as
\[
\begin{aligned}
 u_0&=(0,0,1),&u_4&=(1,0,1),\\
 u_1&=(1,-1,1),&u_2&=(0,1,1).
\end{aligned}
                                                        \tag{P3K.25.2}
\]
Their relation is \(u_0+u_4=u_1+u_2\). This is exactly the Cox
square of P3K.17 under
\[
 B=\begin{pmatrix}1&0\\-1&1\end{pmatrix}.
\]
The auxiliary weights are \((1,1,-1,-1)\), and the physical
delta weights are \((-1,0,0,1)\): their image in the height-zero
lattice is \(u_2-u_0=e_2\). Hence every nonzero two-by-two
minor of the two-circle weight matrix is \(1\) or \(-1\).
The seven-dimensional real master space has the same semifree
residual circle, fixed normal pairs, and inner \(\mathbb{CP}^3\)
face as P3K.17.

The torus coordinates in this square are the literal functions
\[
                  Z_{\mathrm{loc}}=Z,\qquad
                  V_{\mathrm{loc}}=ZV.                 \tag{P3K.25.3}
\]
They are used on the common dense torus to name the Cox maps;
they do not change the global source line. Write the actual
whole-I1 periods in its multidegree-zero source frame as
\[
\begin{aligned}
 h:(Z,V)&\longmapsto(qZ,\alpha(q)V),\\
 \gamma:(Z,V)&\longmapsto(g(q)Z,q\lambda\beta(q)V),
\end{aligned}
                                                        \tag{P3K.25.4}
\]
where the nonzero units have parameter-independent logarithms.
These are the periods of P3K.22.3, with
\(g=\exp(2\pi i\Gamma)\), \(\beta=\exp(2\pi iC)\).
In (P3K.25.3) they are exactly
\[
\begin{aligned}
 h:(Z_{\mathrm{loc}},V_{\mathrm{loc}})
       &\longmapsto(qZ_{\mathrm{loc}},q\alpha V_{\mathrm{loc}}),\\
 \gamma:(Z_{\mathrm{loc}},V_{\mathrm{loc}})
       &\longmapsto(gZ_{\mathrm{loc}},
                          qg\beta\lambda V_{\mathrm{loc}}).
\end{aligned}
                                                        \tag{P3K.25.5}
\]
Their integer labels are \((1,1)\) and \((0,1)\).
They generate all of \(\mathbb Z^2\). Thus the disjoint periodic
collection of square-diagonal tubes is one free period orbit,
with no stabilizer on a chosen master-chart representative.
The quotient of the disjoint collection is that single chart.
It has no residual source or gamma circle.

The original fan uses the \(u_0u_4\) diagonal, and the new fan
uses \(u_1u_2\). These are the two moment faces of P3K.17 in
the reverse order from that statement. All its actual quotient
maps on the common rank-one matrix locus remain the identity
maps in the named global torus coordinates.

\par\medskip\noindent\textbf{The actual scalar loop and the common endpoint gauge.}\enspace

The positive parameter deck on the ordinary cusp collar is
the actual line multiplication \(q^{-1}\). It is the deck
derived from the directed cap-to-ordinary gluing in P3K.10
and P3K.22. Returning its shifted square representative using
the unique gamma period cancels the integral \(q\)-translation.
The remaining bounded torus map in (P3K.25.3) is
\[
       (Z_{\mathrm{loc}},V_{\mathrm{loc}})
           \longmapsto(gZ_{\mathrm{loc}},
                             g\beta\lambda V_{\mathrm{loc}}).
                                                        \tag{P3K.25.6}
\]
Its Cox lift is the literal formula
\[
\begin{aligned}
 z_0&\longmapsto(g^2\beta\lambda)^{-1}z_0,
 & z_4&\longmapsto z_4,\\
 z_1&\longmapsto gz_1,
 & z_2&\longmapsto g\beta\lambda z_2 .
\end{aligned}                                          \tag{P3K.25.7}
\]
The product of the four multipliers is one, so the base
monomial is fixed. The parameter-dependent part has weights
\((-1,0,0,1)\); it is one global physical delta multiplication.
The same formula acts at both fixed charts of the master
space. Its remaining unit multipliers are independent of the
parameter and have the stated logarithms.

This comparison also retains the actual extending origins,
rather than choosing a new loop from (P3K.25.7).
With \(s=\log q/(2\pi i)\), \(v=\log\lambda/(2\pi i)\),
they are
\[
 o=(-\Gamma s,-Cs-s^2/2-sv),\qquad
 \sigma=(\Gamma v,Cv+(v^2-v)/2).
\]
Their difference changes by
\[
 (\sigma-o)(v+\theta)-(\sigma-o)(v)
     =\theta\gamma+\frac{\theta(\theta-1)}2\delta.
                                                        \tag{P3K.25.8}
\]
The first term is the one positive gamma traversal of the
actual family. The second is one common delta gauge. On the
auxiliary source cover \(q=Q^2\), the source deck changes this
angular expression by \(2\theta\delta\). This is likewise the
same physical circle scalar on every master chart and both
cap endpoints. The half-integral one-step ordinary-origin
term is retained exactly as in P3K.10; no false one-step
descent is imposed.

To pass through the real moment bordism, follow each bounded
Cox unit map by its unique positive auxiliary rescaling back
to the prescribed moment value. The derivative of the moment
equation is positive away from the removed origin. Use the
images of the inner and outer truncation faces under these
equivariant isotopies, exactly as in P3K.17. Contract the
parameter-independent units through their chosen logarithms,
with the exterior cutoff and transported collars. The
remaining angular map is physical delta and therefore acts
trivially on the circle orbit and on its oriented normal
three-plane bundle.

\par\medskip\noindent\textbf{All extra boundaries and their actual lines.}\enspace

The inner face has orbit \(S^2_A*S^2_B\). Its free part carries
the actual circle line
\[
                       x=h_A-h_B.
                                                        \tag{P3K.25.9}
\]
The two charged cap maps are those of P3K.17:
\((A,B)=(P,\iota Q)\) at the first end and
\((A,B)=(Q,\iota P)\) at the second. Both pull the line to
\(p+q\), by the explicit orthogonal-complement line
isomorphism. The resulting closed five-dimensional correction
is \(T_g(S^2\times S^2)\), \(g(A,B)=(\iota B,\iota A)\).
Here this \(g\) denotes the conjugate-swap map, not the
source-period unit in (P3K.25.4).

On the two overlap components of this joint cap, (P3K.25.7)
and (P3K.25.8) give transition functions
\(c(\theta)\varphi_{1,0}\) and
\(c(\theta)\varphi_{2,0}\), with the identical scalar
\(c(\theta)\). One change of cap frame by \(c(\theta)^{-1}\)
removes both angular factors. Thus its line and tangent
Pontryagin class pull back from the five-dimensional
correction. Its two six-dimensional characteristic numbers
\(\int x^3,\int p_1x\) vanish.

The fixed locus of the master space is the same actual
three-dimensional cobordism
\[
 (S^2\times I)\setminus
                   (\mathring D^3\sqcup\mathring D^3).
\]
It is obtained by keeping the annulus in the original fixed
curve and the two Cox fixed charts. Its four boundary
line charges are \(0,0,+h_A,-h_B\).
The two zero endpoint charges follow from the literal
\(\mathcal O(-1)\oplus\mathcal O(-1)\) normal charts with
weights \((-1,+1)\), independently of the source line \(M\).
Conjugating the negative-weight summand gives the oriented
Hopf normal model with determinant class
\(-h+h=0\); the adjacent ray relation is unchanged by
every shear fixing \(e_2\).
The original endpoint conifold caps and the joint charged
cap close it to the oriented \(S^2\)-mapping torus of P3K.17.
Its parameter circle is a separate factor: the parameter
gauge is literally physical delta, and the common scalar was
removed by the preceding single frame change. Consequently
the class \(b\) of the additional normal-sphere boundary
satisfies \(b^2=0\), and its base tangent Pontryagin number
is zero. Retaining that entire sphere boundary gives
\[
 \frac{4\int x^3-\int p_1x}{24}
       =\frac{3\int b^2-\int p_1(T\bar X)}{24}=0.
                                                        \tag{P3K.25.10}
\]
This uses no value for the normal bundle's \(p_1\), and no
spin extension across the comparison bordism.

The two actual endpoint conifold mapping tori have opposite
boundary signs in this oriented bordism carrying their
actual Euler lines. The other two boundary contributions
were just proved zero. Hence their characteristic expressions
\((4\int x^3-\int p_1x)/24\), with their prescribed endpoint
spin structures, are equal. P3F.22 then identifies their
oriented seven-sphere return classes.

Finally repeat this actual square construction after the
integral shears
\(\left(\begin{smallmatrix}1&0\\-j&1\end{smallmatrix}\right)\).
They fix delta; the source-period label changes to
\((1,j+1)\), and the gamma label remains \((0,1)\).
The local source and line units change, but they remain
parameter-independent. Formula (P3K.25.7) still has exactly
the same physical-delta parameter factor. The inverse
comparison handles negative \(j\). This proves the stated
equality for every fan (P3K.25.1).
\end{proof}
\subsection{Scope and audit record}

This statement concerns the explicit whole-I1 scalar families
of P3K.10 and P3K.22 and their complete line of delta-edge fans.
It assigns no Jin integer to a fan without a fixed delta curve
and no invariant to a further finite-isotropy flop. The local
repaired-page calculation in
\path{notes/part3_nodal_page_repair.md} remains an exact local
calculation; its global trace comparison is not inferred from
its characteristic fractions.

The matrix certificate is
\path{code/part3_fourfold/whole_i1_shear.py} and its JSON output.
Root independently audited the literal ray and coordinate map,
the complete period lattice, the common physical-delta Cox
parameter character, both actual origins and the auxiliary
source-root gauge, and both charged boundary corrections.
The root audit passed before this proof was promoted.
